% Source: https://artofproblemsolving.com/community/c3h2085967_may_match_school_math_with_advanced_graphing_10_questions
% Stable identity: aops-topic-2085967-q03
On the trip from \(A\) to \(B\) the speed never exceeds \(22\), so it is never \(58\).

On the trip from \(B\) to \(C\) the speed starts at \(0\) (the car is stopped at \(B\)), reaches \(60\) at some moment \(t_0\), and returns to \(0\) at \(C\). Since the speed is continuous, the intermediate value theorem gives a moment before \(t_0\) and a moment after \(t_0\) at which the speed is exactly \(58\). So there are at least \(2\) such moments.

Two is achievable: accelerate steadily from \(0\) to \(60\) and then decelerate steadily to \(0\); the speed passes \(58\) once on the way up and once on the way down.

The least number is \(\boxed{2}\).
